\chapter{Successor and Functional Graphs}
\chapterepigraph{When every node has exactly one outgoing edge, the graph
is a \emph{functional graph}: following the arrows always leads into a
cycle, with trees of nodes feeding into it. This rigid structure -- ``rho
shaped'' tails leading to loops -- makes far-future queries answerable by
binary lifting and cycle lengths computable in one pass.}

\section{One Outgoing Edge Changes Everything}
\label{sec:cses-graph-functional-intro}

In a successor graph each node $x$ has a single successor
$\textit{succ}(x)$. Repeatedly applying $\textit{succ}$ from any start must
eventually repeat a node and enter a cycle, since the state space is
finite.

\begin{patternbox}[Functional-Graph Techniques]
\begin{itemize}
  \item \textbf{``Where am I after $k$ steps?'' for huge $k$} --
        \textbf{binary lifting}: precompute $\textit{succ}^{2^j}$ so any
        $k$ is composed from its binary digits: Planets Queries I.
  \item \textbf{Distance between two nodes on the functional graph} --
        binary lifting plus cycle analysis (a directed analogue of LCA):
        Planets Queries II.
  \item \textbf{Length of the cycle each node eventually enters} --
        one pass marking tails and loops (or Floyd/Brent cycle finding):
        Planets Cycles.
  \end{itemize}
\end{patternbox}

\begin{ideabox}[Binary Lifting Answers $k$ Steps in $O(\log k)$]
Store, for each node, where you land after $2^0,2^1,2^2,\ldots$ steps.
Any $k$-step jump is the composition of the jumps for the set bits of $k$,
so a query costs $O(\log k)$ after $O(n\log k)$ preprocessing -- the same
doubling idea as sparse tables and LCA.
\end{ideabox}

\newpage

\section{Planets Queries I}
\label{sec:cses-planets-queries-i}

\problemstatement{Each of $n$ planets has exactly one teleporter to a
successor planet. Answer $q$ queries: starting from planet $x$ and
teleporting $k$ times (where $k$ can be up to $10^9$), which planet do you
reach?}

\begin{patternbox}
Following one edge per node many times is answered by \textbf{binary
lifting}: precompute where you land after $2^0,2^1,2^2,\ldots$ steps, then
compose the jumps for the set bits of $k$ to answer each query in
$O(\log k)$.
\end{patternbox}

\subsection*{Doubling the jump}

Let $\textit{up}[j][x]$ be the planet reached after $2^j$ teleports from
$x$. The base is $\textit{up}[0][x]=\textit{succ}(x)$, and doubling gives
$\textit{up}[j][x]=\textit{up}[j-1][\textit{up}[j-1][x]]$ -- jump $2^{j-1}$,
then $2^{j-1}$ more. For a query $(x,k)$, walk through the bits of $k$: for
each set bit $j$, replace $x$ with $\textit{up}[j][x]$. After processing all
bits, $x$ is the destination.

\begin{ideabox}[Any k Is a Sum of Powers of Two]
Because $k$ decomposes into its binary digits, a $k$-step jump is the
composition of the precomputed $2^j$-jumps for the bits set in $k$. Thus
$O(n\log k)$ preprocessing yields $O(\log k)$ per query -- the same doubling
trick as sparse tables and LCA.
\end{ideabox}

\subsection*{Algorithm}

\begin{enumerate}
  \item Read $\textit{succ}$; build $\textit{up}[0]$; fill
        $\textit{up}[j]$ by composing the previous level.
  \item For each $(x,k)$: for every set bit $j$ of $k$, advance $x=
        \textit{up}[j][x]$.
  \item Output the final $x$.
\end{enumerate}

\subsection*{Dry run}

\dryrun{$\textit{succ}=[2,3,1]$ (a 3-cycle); query $x=1,k=4$.}

\begin{itemize}
  \item $k=4=100_2$: jump $2^2=4$ steps. $\textit{up}[2][1]=$ four steps
        around the 3-cycle $=1\to2\to3\to1\to2$, landing on $2$.
  \item Answer $2$.
\end{itemize}

\subsection*{C++ implementation}

\begin{lstlisting}
#include <bits/stdc++.h>
using namespace std;
const int LOG = 30;                              // 2^30 > 1e9

int main() {
    int n, q;
    cin >> n >> q;
    vector<vector<int>> up(LOG, vector<int>(n + 1));
    for (int i = 1; i <= n; i++) cin >> up[0][i];  // successor
    for (int j = 1; j < LOG; j++)
        for (int i = 1; i <= n; i++)
            up[j][i] = up[j - 1][up[j - 1][i]];

    while (q--) {
        int x, k;
        cin >> x >> k;
        for (int j = 0; j < LOG; j++)
            if (k & (1 << j)) x = up[j][x];        // apply 2^j-jump for set bits
        cout << x << "\n";
    }
    return 0;
}
\end{lstlisting}

\begin{complexitybox}
\textbf{Time Complexity.} $O(n\log k)$ preprocessing, $O(\log k)$ per query.
\textbf{Space Complexity.} $O(n\log k)$ for the lifting table.
\end{complexitybox}

\begin{takeawaybox}
Binary lifting answers ``$k$ steps ahead'' in a functional graph in
logarithmic time by composing power-of-two jumps. The same table underlies
$k$-th ancestor queries and LCA, and powers the harder Planets Queries II
next.
\end{takeawaybox}

\section{Planets Queries II}
\label{sec:cses-planets-queries-ii}

\problemstatement{Each of $n$ planets has one teleporter. Answer $q$
queries: the minimum number of teleports to go from planet $a$ to planet
$b$, or $-1$ if impossible.}

\begin{patternbox}
A functional graph is ``rho-shaped'': trees of tail nodes feeding into
cycles. Decompose it into cycles and the depth of each node to its cycle,
then a query splits into a \textbf{tail case} (answered by binary lifting)
and a \textbf{cycle case} (answered by modular arithmetic on the cycle).
\end{patternbox}

\subsection*{Tails into cycles}

Peel nodes of in-degree $0$ repeatedly; whatever remains lies on cycles.
Label each cycle, giving its nodes positions $0,1,\ldots$ and recording the
cycle length. Each tail node gets a $\textit{depth}$ (steps to reach its
cycle) and an $\textit{entry}$ (the cycle node it first lands on). For a
query $(a,b)$:
\begin{itemize}
  \item \textbf{Tail case:} if $b$ lies on the path from $a$ into the
        cycle, then $\textit{depth}[a]\ge\textit{depth}[b]$ and jumping
        $\textit{depth}[a]-\textit{depth}[b]$ steps (binary lifting) from
        $a$ lands on $b$; that jump count is the answer.
  \item \textbf{Cycle case:} if $b$ is on the same cycle $a$ reaches, walk
        $\textit{depth}[a]$ steps to $a$'s entry, then around the cycle:
        $\textit{depth}[a]+((\textit{pos}[b]-\textit{pos}[\textit{entry}[a]])
        \bmod L)$.
\end{itemize}
Take whichever applies; otherwise $-1$.

\begin{ideabox}[Split the Query by Where b Lives]
The rho structure means $b$ is reachable from $a$ only if it is on $a$'s
tail before the cycle, or on the cycle $a$ enters. Binary lifting verifies
the tail case exactly; modular position arithmetic handles the cycle case.
\end{ideabox}

\subsection*{Algorithm}

\begin{enumerate}
  \item Build binary lifting for the successor function.
  \item Peel to find cycle nodes; assign cycle ids, positions, lengths;
        compute $\textit{depth}$ and $\textit{entry}$ for tail nodes.
  \item Per query, evaluate the tail and cycle cases; output the minimum or
        $-1$.
\end{enumerate}

\subsection*{Dry run}

\dryrun{$\textit{succ}=[2,3,1,3]$: $1\!\to\!2\!\to\!3\!\to\!1$ cycle, $4\!
\to\!3$ tail.}

\begin{itemize}
  \item Query $(4,1)$: $4$ enters the cycle at $3$ in $1$ step, then $3\to
        1$ is $2$ more steps around the cycle. Answer $3$.
  \item Query $(1,4)$: $4$ is a tail node not reachable from the cycle --
        $-1$.
\end{itemize}

\subsection*{C++ implementation}

\begin{lstlisting}
#include <bits/stdc++.h>
using namespace std;
const int LOG = 20;

int main() {
    int n, q;
    cin >> n >> q;
    vector<int> succ(n + 1);
    vector<vector<int>> up(LOG, vector<int>(n + 1));
    vector<int> indeg(n + 1, 0);
    for (int i = 1; i <= n; i++) { cin >> succ[i]; up[0][i] = succ[i]; indeg[succ[i]]++; }
    for (int j = 1; j < LOG; j++)
        for (int i = 1; i <= n; i++) up[j][i] = up[j - 1][up[j - 1][i]];

    auto jump = [&](int x, int k) {
        for (int j = 0; j < LOG; j++) if (k & (1 << j)) x = up[j][x];
        return x;
    };

    // peel tree nodes; remaining are cycle nodes
    vector<int> deg = indeg, order;
    queue<int> pq;
    for (int i = 1; i <= n; i++) if (deg[i] == 0) pq.push(i);
    vector<bool> removed(n + 1, false);
    while (!pq.empty()) {
        int u = pq.front(); pq.pop();
        removed[u] = true; order.push_back(u);
        if (--deg[succ[u]] == 0) pq.push(succ[u]);
    }

    vector<int> depth(n + 1, 0), entry(n + 1, 0), lead(n + 1, 0), pos(n + 1, 0);
    vector<bool> isCyc(n + 1, false);
    vector<int> cycLen(n + 1, 0);
    int cid = 0;
    for (int i = 1; i <= n; i++) {
        if (removed[i] || isCyc[i]) continue;
        cid++; int c = i, len = 0;
        for (int x = i; !isCyc[x]; x = succ[x]) {
            isCyc[x] = true; pos[x] = len++; lead[x] = cid;
            entry[x] = x; depth[x] = 0;
        }
        cycLen[cid] = len;
    }
    for (int idx = (int)order.size() - 1; idx >= 0; idx--) {   // reverse peel order
        int u = order[idx];
        depth[u] = depth[succ[u]] + 1;
        entry[u] = entry[succ[u]];
        lead[u] = lead[succ[u]];
    }

    while (q--) {
        int a, b;
        cin >> a >> b;
        long long ans = -1;
        if (a == b) { cout << 0 << "\n"; continue; }
        // tail case: b on a's path into the cycle
        if (depth[a] >= depth[b]) {
            int steps = depth[a] - depth[b];
            if (jump(a, steps) == b) ans = steps;
        }
        // cycle case: b on the same cycle a reaches
        if (isCyc[b] && lead[b] == lead[a]) {
            int L = cycLen[lead[a]];
            long long steps = depth[a] + ((pos[b] - pos[entry[a]]) % L + L) % L;
            if (ans == -1 || steps < ans) ans = steps;
        }
        cout << ans << "\n";
    }
    return 0;
}
\end{lstlisting}

\begin{complexitybox}
\textbf{Time Complexity.} $O(n\log n)$ preprocessing, $O(\log n)$ per query.
\textbf{Space Complexity.} $O(n\log n)$ for the lifting table.
\end{complexitybox}

\begin{takeawaybox}
Distances in a functional graph split by structure: binary lifting settles
the tree-tail case, modular cycle arithmetic settles the cycle case. Seeing
the rho shape -- tails feeding cycles -- is the key to organising the query
logic.
\end{takeawaybox}

\section{Planets Cycles}
\label{sec:cses-planets-cycles}

\problemstatement{Each of $n$ planets has one teleporter. For each starting
planet, output how many planets you visit before returning to a planet you
have already seen (i.e.\ the number of teleports until the first repeat).}

\begin{patternbox}
In a functional graph, starting anywhere you walk a tail and then loop a
cycle. The count until the first repeat is $\textit{depth}+\textit{cycle
length}$: the tail steps to reach the cycle plus the cycle's size. So the
same rho decomposition as Planets Queries II answers every node at once.
\end{patternbox}

\subsection*{Tail length plus cycle length}

Peel to identify cycle nodes; for each cycle record its length. A node on a
cycle simply revisits after $\textit{cycleLen}$ steps. A tail node walks
$\textit{depth}$ steps to reach the cycle, then $\textit{cycleLen}$ more to
return to its entry point -- so its answer is
$\textit{depth}+\textit{cycleLen}(\textit{entry})$. Computing $\textit{depth}$
and the cycle each node leads to (in reverse peel order) gives all answers in
one pass.

\begin{ideabox}[Every Walk Is Tail Then Loop]
The first repeat is always the cycle entry, reached again after one full
loop. Hence the visited count is exactly tail length plus cycle length,
independent of where on the cycle you enter.
\end{ideabox}

\subsection*{Algorithm}

\begin{enumerate}
  \item Peel tree nodes; label cycles with their lengths.
  \item Cycle node answer $=\textit{cycleLen}$; tail node answer
        $=\textit{depth}+\textit{cycleLen}(\textit{lead})$.
  \item Output every node's answer.
\end{enumerate}

\subsection*{Dry run}

\dryrun{$\textit{succ}=[2,3,1,3]$: cycle $1\!\to\!2\!\to\!3$ (length $3$),
tail $4\!\to\!3$.}

\begin{itemize}
  \item Nodes $1,2,3$: answer $3$ each. Node $4$: depth $1$ + cycle $3$
        $=4$.
  \item Output $3\ 3\ 3\ 4$.
\end{itemize}

\subsection*{C++ implementation}

\begin{lstlisting}
#include <bits/stdc++.h>
using namespace std;

int main() {
    int n;
    cin >> n;
    vector<int> succ(n + 1), indeg(n + 1, 0);
    for (int i = 1; i <= n; i++) { cin >> succ[i]; indeg[succ[i]]++; }

    vector<int> deg = indeg, order;
    queue<int> pq;
    for (int i = 1; i <= n; i++) if (deg[i] == 0) pq.push(i);
    vector<bool> removed(n + 1, false);
    while (!pq.empty()) {
        int u = pq.front(); pq.pop();
        removed[u] = true; order.push_back(u);
        if (--deg[succ[u]] == 0) pq.push(succ[u]);
    }

    vector<int> depth(n + 1, 0), lead(n + 1, 0), cycLen(n + 1, 0);
    vector<bool> isCyc(n + 1, false);
    int cid = 0;
    for (int i = 1; i <= n; i++) {
        if (removed[i] || isCyc[i]) continue;
        cid++; int len = 0;
        for (int x = i; !isCyc[x]; x = succ[x]) { isCyc[x] = true; lead[x] = cid; len++; }
        cycLen[cid] = len;
    }
    for (int idx = (int)order.size() - 1; idx >= 0; idx--) {
        int u = order[idx];
        depth[u] = depth[succ[u]] + 1;
        lead[u] = lead[succ[u]];
    }

    for (int i = 1; i <= n; i++)
        cout << depth[i] + cycLen[lead[i]] << " ";
    cout << "\n";
    return 0;
}
\end{lstlisting}

\begin{complexitybox}
\textbf{Time Complexity.} $O(n)$ -- peeling plus one reverse pass.
\textbf{Space Complexity.} $O(n)$.
\end{complexitybox}

\begin{takeawaybox}
Steps-until-repeat in a functional graph is tail depth plus cycle length,
computed for every node from one rho decomposition. Recognising the tail-
then-loop structure turns a per-node simulation into a single linear pass.
\end{takeawaybox}

